\section{引言}
\label{sec:introduction}

生成视频检测需要应对持续变化的生成模型。
免训练方法复用预训练表示或生成模型的统计特性，避免为每类新生成器训练真假分类器。
近期工作分别从特征变换的稳定性~\cite{choi2025warpad}、
帧间变化~\cite{zheng2025d3,brokman2026overcoherence}及真实参考似然~\cite{benhayun2026stall}
构造检测分数。
这使研究重点从扩大分类器转向一个更具体的问题：冻结表示中的哪些信息，应当以何种方式被测量？

时间线索既可以体现为变化，也可以体现为过度连贯。
D3~\cite{zheng2025d3}度量相邻特征距离的二阶波动，
Over-Coherence~\cite{brokman2026overcoherence}利用局部或整体偏高的帧间相似性，
SPLIT~\cite{hyun2026split}进一步比较图块路径粗糙度与邻接位移。
STALL~\cite{benhayun2026stall}则用真实统计联合解释全局外观和归一化一阶变化。
这些方法确立了时间与局部证据的价值，但尚未替代对
\emph{局部高维二阶方向在真实度量下的增量}的直接检验。

本文的出发点是测量与聚合的顺序。
不同位置的变化可能在空间平均中抵消；具有相同长度的残差，也可能落在真实统计中具有不同代价的方向。
因此，我们在空间汇总前保留各位置的归一化二阶残差，
让真实协方差定义方向度量，再聚合位置评分。
这一设计利用了局部动态分布的信息，同时保留与全局证据共享视觉编码的简单实现。

我们将该方法称为\method。
固定真实参考后，同一套评分规则用于该域的各个生成器；推理无需生成器身份。
局部方向表示、真实统计度量和固定预算片段观察共同构成检测流程，
其中局部方向的补充价值通过匹配对照验证。
本文的贡献为：
\begin{itemize}
  \item \textbf{局部方向的真实统计评分。}
  在冻结表示上先测量逐位置二阶方向、后汇总证据，并解释该顺序保留的方向离散信息。
  \item \textbf{区分表示收益与额外参考支持。}
  在相同真实预算与观察条件下，通过聚合、阶次和位置数控制检验局部增量，并以新真实拟合池验证其重复性。
\end{itemize}
